% !TeX root = ../main.tex
Quantile trend filtering combines the check loss with an $\ell_1$ penalty on discrete differences of the trend. We develop Bayesian-PMALA, a posterior sampler for quantile trend filtering that learns the regularization strength and the likelihood scale jointly with the trend. Augmenting the asymmetric Laplace likelihood with its normal--exponential representation makes the likelihood Gaussian in the trend. Under a global Laplace prior on trend differences, the whole negative log-conditional of the trend is then a weighted quadratic plus an $\ell_1$ difference penalty. The trend can therefore be updated by an exact proximal Metropolis-adjusted Langevin (P-MALA) step. The Moreau--Yosida envelope is built on the entire conditional posterior, through a weighted trend filtering problem, rather than on a smoothed check loss, and the Metropolis--Hastings correction uses the exact density. The latent variables, the penalty and the scale have conjugate updates. Every linear system is banded, and the proximal map is computed by ADMM with a factorization shared by the forward and reverse maps, so the work per sweep grows linearly with the sample size for a fixed inner iteration budget. For several quantile levels, a joint version samples each trend between its neighbours. Only the proximal map changes, becoming a box-constrained problem solved with banded interior-point steps, and no retained draw crosses.

In simulations with varying-smoothness and smooth signals under Gaussian, Beta, and mixed-normal noise, Bayesian-PMALA attains accuracy comparable to a locally adaptive horseshoe Gibbs sampler with a simpler hierarchy and banded computations, and accuracy comparable to ADMM-based quantile trend filtering while providing posterior inference on the trend, the penalty and the scale. In a peak-detection design for air-quality sensors, it recovers the peak-free baseline far more accurately than the horseshoe sampler and comes close to the ADMM benchmark in baseline recovery and detection, and an application to fenceline sensor data illustrates the method at $n = 8{,}640$.

% Previous abstract (kept for reference, 2026-09-26):
% Quantile trend filtering combines the check loss with an $\ell_1$ difference penalty. We develop a Bayesian extension of its ADMM-based optimization formulation, learning the regularization strength and likelihood scale jointly with the trend. A three-block Gibbs scheme combines conjugate scalar updates with a proximal Metropolis-adjusted Langevin trend update, using ADMM to solve the proximal subproblems. The hierarchy uses one global penalty per quantile rather than a vector of random local shrinkage scales. The constrained proximal system permits reuse of a banded factorization and linear work in the sample size for a fixed inner iteration budget.
%
% Joint estimation incorporates non-crossing constraints through box projections and sequential quantile updates. Auxiliary sentinel quantiles provide dynamic bounds for extreme target levels. Simulations use varying-smoothness and smooth signals under Gaussian, Beta, and mixed-normal noise. At interior levels, the sampler has the lower error in the Gaussian and mixed-normal varying-smoothness cells while the horseshoe is decisively better under Beta noise, and it improves on both Bayesian competitors in every smooth-signal configuration. Its smooth-signal errors are of a similar magnitude to the available ADMM results. A further experiment reproduces a published peak-detection design, in which a low-quantile baseline is estimated from a series contaminated by transient peaks: there the locally adaptive competitor follows the peaks and is two to fourteen times less accurate on the baseline, while the global penalty detects the peaks as well as the optimization benchmark without selecting a regularization parameter.
%
% The method provides a compromise between efficient point optimization and the more elaborate locally adaptive Gibbs hierarchy: it adds posterior sampling to the former and simplifies the latter's computational structure. Incomplete ADMM replications and limitations of the reported intervals qualify the comparison; posterior output alone does not establish calibration.
